\section{DQN Algorithm}

\begin{frame}{DQN Algorithm}
    \begin{center}
    \footnotesize
    \begin{algorithm}[H]
        \DontPrintSemicolon
        \caption{Deep Q-Network (DQN) with Experience Replay}
        Initialize replay memory $\mathcal{D}$ to capacity $N$\;
        Initialize action-value network $Q$ with random weights $\theta$\;
        \For{$\text{episode} = 1$ \KwTo $M$}{
            Observe initial state $s$\;
            \For{$t = 1$ \KwTo $T$}{
                With probability $\varepsilon$ select a random action $a$; otherwise $a = \arg\max_{a'} Q(s, a'; \theta)$\;
                Execute $a$; observe reward $r$ and next state $s'$\;
                Store transition $(s, a, r, s')$ in $\mathcal{D}$;\quad $s \leftarrow s'$\;
                Sample a random minibatch of transitions $(s_i, a_i, r_i, s_i')$ from $\mathcal{D}$\;
                Set $y_i = \begin{cases} r_i & \text{if } s_i' \text{ is terminal} \\ r_i + \gamma \max_{a'} Q(s_i', a'; \theta) & \text{otherwise} \end{cases}$\;
                Take a gradient descent step on $\big(y_i - Q(s_i, a_i; \theta)\big)^2$ w.r.t.\ $\theta$\;
            }
        }
    \end{algorithm}
    \end{center}
    \footnotetext{$\varepsilon$-greedy exploration---covered in Day\,1. $\varepsilon$ is annealed from high to low over training. Prediction and target share the \emph{same} weights $\theta$---the instability this causes motivates the next slide.}
\end{frame}

\begin{frame}{DQN with Target Network}
    \begin{itemize}
        \item DQN, as stated above, will have learning problems because the target and the prediction are not independant as they both rely on the same network.
        \item It's like a dog chasing it's own tail.
        \pause
        \item \textbf{Solution:} Use two separate $Q$-value estimators, each of which is used to update the other.
        \item The target values are calculated using a target Q-network. The target Q-network's parameters are updated to the current networks every $C$ time steps.
        \item Target network prevents the network from spiraling around.
    \end{itemize}
\end{frame}

\begin{frame}{DQN with Target Network Algorithm}
    \begin{center}
    \footnotesize
    \begin{algorithm}[H]
        \DontPrintSemicolon
        \caption{DQN with a Target Network}
        Initialize replay memory $\mathcal{D}$ to capacity $N$\;
        Initialize online network $Q$ with random weights $\theta$\;
        Initialize target network $\hat{Q}$ with weights $\theta^- \leftarrow \theta$\;
        \For{$\text{episode} = 1$ \KwTo $M$}{
            Observe initial state $s$\;
            \For{$t = 1$ \KwTo $T$}{
                With probability $\varepsilon$ select a random action $a$; otherwise $a = \arg\max_{a'} Q(s, a'; \theta)$\;
                Execute $a$; observe reward $r$ and next state $s'$\;
                Store transition $(s, a, r, s')$ in $\mathcal{D}$;\quad $s \leftarrow s'$\;
                Sample a random minibatch of transitions $(s_i, a_i, r_i, s_i')$ from $\mathcal{D}$\;
                Set $y_i = \begin{cases} r_i & \text{if } s_i' \text{ is terminal} \\ r_i + \gamma \max_{a'} \hat{Q}(s_i', a'; \theta^-) & \text{otherwise} \end{cases}$\;
                Take a gradient descent step on $\big(y_i - Q(s_i, a_i; \theta)\big)^2$ w.r.t.\ $\theta$\;
                Every $C$ steps set $\theta^- \leftarrow \theta$\;
            }
        }
    \end{algorithm}
    \end{center}
    \footnotetext{Mnih et al.\ (2015), \href{https://www.nature.com/articles/nature14236}{Human-level control through deep reinforcement learning}, \emph{Nature}. The target $y_i$ now uses the frozen network $\hat{Q}(\cdot;\theta^-)$.}
\end{frame}